\documentclass[10pt]{article}

\usepackage[T1]{fontenc}
\usepackage{lmodern}
\usepackage[margin=0.8in]{geometry}
\usepackage[hidelinks]{hyperref}
\hypersetup{
  pdfauthor={},
  pdftitle={Cover Letter - Computers and Graphics}
}
\setlength{\parindent}{0pt}
\setlength{\parskip}{4pt}

\begin{document}

\begin{flushright}
September 23, 2026
\end{flushright}

Dear Editor-in-Chief and Editors of \emph{Computers \& Graphics},

Please consider our Research Paper,
\begin{center}
\emph{Dynamic Prototype Routing for Lightweight Image Super-Resolution},
\end{center}
for publication in \emph{Computers \& Graphics}.

Single-image super-resolution is a picture-manipulation problem with direct
relevance to visual computing, display rendering, and resource-constrained image
applications. The manuscript studies how semantically related image tokens can
be grouped and processed efficiently rather than enlarging the reconstruction
backbone.

The main contribution is Dynamic Prototype Routing, a drop-in replacement for
the center-based routing in CATANet. It constructs prototypes from the current
image, separates prototype generation from confirmation with a learnable query
bank, retains assignment confidence for token ordering, gating, and fallback,
and uses a bounded temperature with a balancing loss to prevent prototype
collapse. The design is evaluated on five standard benchmarks at three
upscaling factors. Controlled three-seed ablations compare the dynamic
prototypes directly with a reimplemented history-center variant and isolate the
routing components.

\textbf{Closest prior art (published):}

Xin Liu, Jie Liu, Jie Tang, and Gangshan Wu, ``CATANet: Efficient
Content-Aware Token Aggregation for Lightweight Image Super-Resolution,''
\emph{Proceedings of the IEEE/CVF Conference on Computer Vision and Pattern
Recognition (CVPR)}, pp. 17902--17912, 2025.
\url{https://doi.org/10.1109/CVPR52734.2025.01668}.

CATANet maintains routing centers across batches through an exponential moving
average and uses hard token assignments. The submitted work keeps its backbone
and external tensor interface fixed but replaces that routing mechanism with
per-image prototype construction, explicit confidence propagation, decoupled
prototype confirmation, and anti-collapse regularization. DPRNet does not claim
an efficiency advantage over CATANet: it incurs a modest cost increase while
remaining lighter in Multi-Adds than the SwinIR-light and SRFormer-light
baselines under the common input convention. Its contribution is a more
adaptive and interpretable routing mechanism, supported by matched multi-seed
comparisons and routing diagnostics.

The manuscript is original, has not been published previously, and is not under
consideration elsewhere. It is not an extended version of a conference paper.
All authors have approved the manuscript and its submission. The authors declare
no competing interests. All benchmark images are from publicly available
research datasets, and no new data were generated.

The reference list begins on page 13, so the scientific content extends only
slightly beyond the journal's typical 12-page guidance. This small extension is
needed to retain the three-scale benchmark comparison, controlled multi-seed
ablations, visual results, and routing diagnostics that substantiate the method.

The manuscript and this cover letter have been anonymized for double-anonymized
review. Complete author and correspondence information is provided in the
separate title-page file.

Thank you for your consideration.

Sincerely,

The Authors

\end{document}
